% =========================================================================
% §1 Introduction
% Budget: ~1.5 pp
% Drafting order: LATE (after Results, Discussion, Related Work are stable).
%
% Content blueprint (RESEARCH_paper_outline.md §4 / §1):
%   - Coding agents now interleave LLMs with retrieval; retrieval
%     implementations vary widely (agentic grep, file-dependency PageRank,
%     semantic search, structural codebase index).
%   - Open empirical question: inside a fixed coding-agent harness on a fixed
%     model, does adding a structural codebase index actually change cost
%     or resolve?
%   - Scope: open-source harnesses; one model held fixed (Opus 4.7);
%     causal ON/OFF ablation + cross-harness comparator.
%   - Headline preview (multi-seed, three seeds): View B acc@5 ablation
%     +39.6 pp paired Wilcoxon p<0.0001, resolve ablation +7.9 pp paired
%     Wilcoxon p=0.003, lower $/solved with per-cell cost statistically
%     null. Cross-harness vs OpenCode marginal at p=0.087 / p=0.080.
%     [Updated 2026-06-16 after multi-seed analysis.]
%   - Contribution statement: first leak-audited, model-controlled, causal
%     ablation of a shipped structural codebase index in a coding-agent
%     harness.
%
% Constraints:
%   - Closed-source harnesses (Claude Code, Cursor, Windsurf) out of scope by
%     design; called out in one sentence.
%   - No motivating-example opener. (Decision: Option C.)
% =========================================================================

\section{Introduction}
\label{sec:intro}

Coding agents now interleave LLMs with retrieval over the working repository, and retrieval
implementations vary widely across deployed harnesses. The implementations span a
spectrum: agentic grep over the working copy, file-dependency repo maps, semantic and
graph search, and structural codebase indices built once per repository
(\S\ref{sec:related}). Inside a fixed coding-agent harness on a fixed model, does adding a
structural codebase index actually change cost or resolve? We answer the question for
open-source harnesses with the model held fixed (Claude Opus~4.7
\citep{claudeopus47}, \S\ref{sec:design-model}); closed-source harnesses (Claude Code, Cursor, Windsurf) are out
of scope by design. The experimental design isolates the index causally by toggling it on
and off inside one harness while everything else stays identical, and cross-checks the
result against an agentic-grep comparator (\S\ref{sec:design-arms}).

The field has not had a clean answer because controlled measurements are scarce. Most
prior work either compares whole harnesses, where retrieval is confounded with prompt,
tool surface, and control loop, or evaluates retrieval components in isolation against
acc@$k$, without the downstream agentic loop that turns ranking into a fix. The
question whether ``grep is all you need'' has been asked recently in the agentic-search
literature for memory-style document retrieval \citep{grepallyouneed2026}, with grep
favored over vector retrieval; we ask the code-task counterpart, where the candidate
beyond grep is not a vector index but a structural codebase index (semantic + lexical +
call-graph). A structural codebase index is also expensive to build and operate, so if
the resolve gain is small and the cost premium is large, the index does not pay for
itself in a deployment. The integrity bar for benchmark evaluation has risen in
parallel: recent audits documented solution leakage in issue text
\citep{swebenchplus2024}, memorization of in-benchmark repositories
\citep{swebenchillusion2025}, and substantial score inflation from formal issue text
relative to realistic user phrasing \citep{savingswebench2025}, so any positive result
needs to survive a leak audit before it counts.

Three arms ran against the same SWE-PolyBench Verified and SWE-bench Pro public instances
(91 instances; Go, Java, Python) with Claude Opus~4.7 fixed throughout, across three seeds:
SC-ON (the SuperCoder \citep{supercoder2024} harness with the index on), SC-OFF (the same
harness with the two engine tools removed, every other component identical), and OpenCode
\citep{opencode2025} (an agentic-grep comparator). Every cell ran
inside a hardened per-task sandbox with a fail-closed git scrub and a post-run leak audit
(\S\ref{sec:integrity}). On the causal within-harness ablation
(\S\ref{sec:results-ablation}), the index moves View~B acc@5 from 44.3\% to 84.5\% across
seeds (paired Wilcoxon $p<0.0001$) and resolve from 41.9\% to 50.4\% (paired Wilcoxon
$p=0.003$), and yields lower cost per solve with a statistically null per-cell cost
difference. On the cross-harness validity check
(\S\ref{sec:results-crossharness}), SC-ON matches or modestly favors OpenCode on resolve
(50.4\% vs.\ 45.3\% mean, paired Wilcoxon $p=0.087$) and on View~B acc@5
(84.5\% vs.\ 75.3\% mean, paired Wilcoxon $p=0.080$) at no cost penalty. The
structural codebase index does not duplicate behavior that competent agentic grep already
reaches; at minimum, it does not regress the agent.

We report the first leak-audited, model-controlled, causal ablation of a shipped
structural codebase index inside a coding-agent harness, paired with a cross-harness
validity check against an agentic-grep comparator. The result reframes the deployment
question from ``is a structural codebase index too expensive to run alongside agentic
grep'' (on these benchmarks, the answer is no: \$2.30 mean across seeds against
OpenCode's \$2.92, favorable on \$/solved) to ``does the workload include multi-file
changes where structural ranking pays off'' (\S\ref{sec:discussion}). Released alongside
the paper are the per-cell exclusion ledger (\S\ref{sec:integrity-ledger}), the leak-audit
script, the dual-view localization extractor, and the full results database, so every
number in \S\ref{sec:results} is reproducible from the released artifacts.
